\documentclass[tikz,border=2mm]{standalone}

% Load the new OT circuits package
\usepackage{otcircuits}

\begin{document}
\begin{otcircuit}[scale=1.2]
    % Standard RC Circuit in OTBlue (default)
    \draw (0,0) to[V=$V_s$] (0,3)
          to[R=$R$, -*] (3,3)
          to[C=$C$] (3,0)
          to[short] (0,0);

    % Ground in OTGreen (safe)
    \draw (3,0) to[short, -*] (3,-0.5) node[ground, ot-safe]{};

    % Differentiation: Highlighting a specific path or component
    % Feedback loop example with Op-Amp
    \begin{scope}[xshift=6cm]
        \node[op amp] (oa) at (0,0) {};

        % Inputs in OTPurple (signal)
        \draw (oa.-) to[short, -*] ++(-1,0) node[left, ot-signal]{$V_{in}$};
        \draw (oa.+) to[short] ++(-0.5,0) node[ground, ot-safe]{};

        % Feedback path in OTOrange (highlight)
        \draw (oa.-) to[short, *-] ++(0,1.5) coordinate(fb)
              to[R=$R_f$, ot-highlight] (fb -| oa.out)
              to[short, -*] (oa.out);

        \draw (oa.out) to[short, -o] ++(1,0) node[right]{$V_{out}$};

        % Label in OTDark (explicit)
        \node[color=OTDark] at (0,-1.5) {Inverting Amplifier Schematic};
    \end{scope}
\end{otcircuit}
\end{document}
